\clearpage
\fancyhead[L]{\small CONNAISSANCE DES TEMS.}
\hypertarget{native-650-lower}{}
\expandafter\def\csname HistoricalNoteText37\endcsname{Gegenwärtige Anzeige ist im August v. J. geschrieben.}
\begin{center}\small Göttingische gelehrte Anzeigen. Stück 6. Seite 49..56. 1828 Januar 10.\\[3pt]\rule{.72\linewidth}{.3pt}\end{center}

\emph{Connaissance des tems, ou des mouvemens célestes, à l'usage des astronomes et des navigateurs, pour l'an \textup{1829}. Publiée par le bureau des longitudes.} Paris 1826. Bei Bachelier. 382 S. gr. Octav. 

Die Einrichtung des astronomischen Kalenders ist noch unverändert dieselbe, wie in den frühern Jahrgängen. Das Verzeichniss der geographischen Längen und Breiten, welches gegen 2000 Oerter enthält, möchte einer strengen Revision oder Umarbeitung bedürftig sein, welche auch nach einer, in den letzten Jahrgängen stehend gewordenen Bemerkung, vom Längenbureau einer besondern Commission aufgetragen ist, und deren Resultate wir in Kurzem zu erwarten haben. --- Die \emph{Zusätze} füllen 163 Seiten und enthalten folgende Artikel. Tafel zur Abkürzung der Berechnung der mit Hülfe des Polarsterns beobachteten Polhöhen und Azimuthe von \textsc{Puissant}. Zur Berechnung dieser nicht unbequemen Hülfstafeln haben die astronomischen Beobachtungen, die an vielen Punkten des trigonometrischen Netzes für die neue Karte von Frankreich angestellt sind, Veranlassung gegeben. Sie beziehen sich auf keine bestimmte Declination, und würden daher auch für Beobachtungen anderer vom Pole nicht zu weit entfernter Sterne, ausser dem Meridian, sich anwenden lassen. --- Bericht an die K. Academie der Wissenschaften über einen Aufsatz des Hrn. \textsc{Puissant}, die Bestimmung der Gestalt der Erde aus geodätischen und astronomischen Messungen betreffend, von \textsc{Legendre} und \textsc{Mathieu}. Der Aufsatz bezieht sich zunächst auf die in Frankreich gemessenen Stücke von zwei Parallelkreisen, und enthält dem Bericht \SourceLeaf{651} zufolge Rechnungsmethoden für die dabei vorkommenden Aufgaben. Da der Aufsatz erst künftig in dem Memorial du dépôt de la guerre gedruckt werden soll, und der Bericht nicht ins Einzelne eingeht, so können wir über den Werth der vom Hrn. \textsc{Puissant} aufgestellten Vorschriften noch nicht urtheilen. Blos in Beziehung auf die Ausgleichung der Winkel eines Dreieckssystems, welches zwei gemessene Grundlinien mit einander verbindet, und wo der berechnete Werth der zweiten Grundlinie etwas verschieden von dem Gemessenen ausfällt, wird Hrn. \textsc{Puissants} Verfahren angeführt. Dieses läuft darauf hinaus, die an den Winkeln eines jeden Dreiecks anzubringenden Correctionen dem Fehler der Summe der drei Winkel proportional zu setzen. Die Berichterstatter bemerken, dass dieses Princip in der Anwendung bequem, aber der Beweis der Rechtmässigkeit noch nicht geführt sei. Wir setzen hinzu, dass die \emph{richtige} Auflösung jener Aufgabe nach den Grundsätzen der Wahrscheinlichkeitsrechnung leicht gefunden, und dadurch die Unzulässigkeit jenes Princips erwiesen wird. --- Ueber die beiden grossen Ungleichheiten in der Bewegung des Jupiter und Saturn von \textsc{Laplace}. Dieser Aufsatz betrifft diejenigen Theile dieser Ungleichheiten, welche vom Quadrate der störenden Kräfte abhängig sind. Nach einer vom Hrn. \textsc{Plana} gemachten Bemerkung ist es unrichtig, zwischen diesem zweiten Theile der grossen Gleichung in der Saturnsbewegung und dem zweiten Theile der grossen Gleichung in der Bewegung des Saturn, dasselbe Verhältniss anzunehmen, welches zwischen den Haupttheilen (die von der ersten Potenz der störenden Masse abhängen) Statt findet. \textsc{Laplace} entwickelt hier die Relation dieser Theile genauer, allein das Resultat ist unvereinbar mit den von \textsc{Plana} gefundenen numerischen Werthen: es scheint also dieser Gegenstand noch weiterer Untersuchung bedürftig zu sein, und es wäre zu wünschen, dass die ganze Theorie der gegenseitigen Störungen des Jupiter und Saturn von Grund aus auf eine von der bisherigen ganz verschiedene Art behandelt würde. --- Bemerkungen über verschiedene Gegenstände der Mécanique Céleste, von demselben, betreffen (zum Theil als Erwiderungen einiger dagegen gemachten Erinnerungen von Hrn. \textsc{Plana}) die Bewegung der Bahn des äussersten Saturnstrabanten, eine periodische Gleichung des Merkur, und die Wirkung der Fixsterne auf das Sonnensystem. --- Ueber die Messung eines Bogens des Parallelkreises in der Breite von $45^\circ$ von \textsc{Brousseaud} und \textsc{Nicollet}. Dieser sehr interessante Aufsatz gibt Nachricht von der Operation, durch welche, vermittelst Pulversignale, der astronomische Längenunterschied zwischen dem Mont-Cenis und dem Thurm von Marennes, an der Mündung der Garonne, in den Jahren 1822 u. 1823 bestimmt worden ist. Der ganze Bogen umfasst 8 Längengrade und die Pulversignale wurden an fünf Zwischenpunkten gegeben. Im Jahre 1822 wurde auf diese Weise der Längenunterschied zwischen dem Mont-Cenis und dem Puy d'Isson im Departement des Puy de Dome unter Zusammenwirken der gedachten französischen Beobachter und der italiänischen Astronomen \textsc{Carlini} und \textsc{Plana} bestimmt; im Jahre 1823, von den französischen Beobachtern allein, das westliche blos auf französischem Gebiete liegende Stück. Der Aufsatz enthält theils alle französischer Seits gemachten astronomischen Beobachtungen, theils die daraus folgenden Resultate für die Längenunterschiede; die letztern sind noch vermittelst der ähnlichen in den Jahren 1821 und 1824 von den italiänischen Astronomen ausgeführten und zum Theil bereits öffentlich bekannt gemachten Operationen bis zur Sternwarte von Padua ausgedehnt. So umfasst diese grosse Unternehmung bereits einen Bogen von 13 Längengraden; später ist derselbe noch bis Fiume erweitert, und es ist Hoffnung zu einer noch weitern Verlängerung vorhanden. Die geodätische Bestimmung des Bogens von Marennes bis Fiume beruht auf 106 Dreiecken erster Ordnung, wovon 90 von französischen, die übrigen von österreichischen und sardinischen Ingenieurs gemessen sind. Es ist sehr zu wünschen, dass auch diese Dreiecke bald ausführlich bekannt gemacht werden, und wir haben dazu gegründete Hoffnung, da nach einer am Schluss des Aufsatzes gemachten Anzeige \emph{alles} auf die Messung dieses Parallelkreisbogens Bezug habende im 9ten Bande des Mémorial du dépôt de la guerre gedruckt werden soll. Ueber die in gegenwärtigem Aufsatze \SourceLeaf{652}enthaltenen Beobachtungen bieten sich noch einige Bemerkungen dar. Die Zeitbestimmung wurde auf absolute mit \textsc{Bordaischen} Repetitionskreisen beobachtete Zenithdistanzen von Sternen gegründet, deren scheinbare Positionen aus \textsc{Schumacher}'s Hülfstafeln entlehnt wurden. Da bei dem in Rede stehenden Geschäft die grösste Genauigkeit der Zeitbestimmung von höchster Wichtigkeit ist, so muss man bedauern, dass statt jenes Verfahrens nicht lieber an den Beobachtungsplätzen transportable Mittagsfernröhre aufgestellt wurden; bei so wichtigen weitumfassenden Operationen hätte die daraus entstehende Vermehrung der Kosten nicht in Betracht kommen können. Bei den Zeitbestimmungen aus Sternen auf verschiedenen Seiten des Meridians finden sich Unterschiede, die über drei Zeitsecunden gehen, wobei es bedenklich bleibt, blos den Mittelwerth zu nehmen. Auch scheint bei einigen Punkten die Anzahl der Abende, wo die Pulversignale beobachtet wurden, zu klein, um über die absolute Zuverlässigkeit der Resultate völlige Beruhigung zu geben. So wurde der Längenunterschied zwischen Solignat und La Jonchere an drei Abenden, jedesmal durch 10 Pulversignale bestimmt, nemlich 

\begin{center}\small\setlength{\tabcolsep}{3pt}
\begin{tabular}{lrll}
August & 12 & . . . . . . . . . . . . . . . . . & $6^{\mathrm m}\ 47^{\mathrm s}86$ \\
 & 19 & . . . . . . . . . . . . . . . . . & $6\quad49.81$ \\
 & $22'$ & . . . . . . . . . . . . . . . . . & $6\quad50.17$ \\
\end{tabular}\end{center}
\noindent Man konnte durchaus keine Ursache ausfindig machen, die beträchtliche Abweichung des ersten Resultats von den beiden andern zu erklären; man schloss jenes ganz aus, und hielt sich nur an das Mittel von diesen. Die Messung des Bogens zwischen La Jonchere und Saint Preuil beruht sogar nur auf den Beobachtungen der Pulversignale einer einzigen Nacht, da der Berg Puy Cogneux, wo letztere gegeben wurden, in St. Preuil, wie es scheint, nur bei stärkerer terrestrischer Refraction sichtbar war. Man war hier bei der Auswahl der Plätze nicht genug auf seiner Hut gewesen; da die grosse Veränderlichkeit der terrestrischen Refraction eine bekannte und bei grossen geodätischen Operationen täglich in die Augen fallende Erscheinung ist, so hätte, ehe St. Preuil zum Beobachtungsplatze gewählt wurde, die Möglichkeit der Gefahr, dass das zwischenliegende hohe Terrain den Signalplatz bei kleinerer oder mittlerer Refraction verdecken könne, allerdings berücksichtigt werden sollen. Zur Entschuldigung dient inzwischen das für so ausgedehnte Operationen zu kleine Personal und die beschränkte Zeit. Immer aber bleibt es auffallend, dass der in St. Preuil angestellte Beobachter anfangs sogar ohne eine \emph{genaue} Kenntniss der Richtung des Signalplatzes war, zu welcher Kenntniss, da Puy Cogneux ein Hauptdreieckspunkt gewesen war, die Mittel doch leicht zu finden sein mussten. Der Werth des Längengrades des Parallelkreises in $45^\circ\,43'\,12''$ Breite wird zuletzt aus dem ganzen Bogen von Marennes bis Padua zu 77865,75 Meter berechnet; den mittlern in diesem Resultate zu befürchtenden Fehler findet Ref. aus der Vergleichung der einzelnen sechs Bestandtheile 25 Meter; allein in Rücksicht der vorhin erwähnten Umstände möchte die zurückbleibende Unzuverlässigkeit eher noch grösser, und daher die Resultate, welche für die Abplattung des Erdsphäroids aus der Vergleichung dieses Längengrades mit verschiedenen gemessenen Breitegraden gezogen werden, noch immer nicht von sehr grossem Gewicht sein. — Ueber die Länge von Port-Praslin in Neu-Irland, von dem Fregatten-Capitain Duperrey. Eine neue Berechnung der von Veron im Jahr 1768 beobachteten Sonnenfinsterniss nach \textsc{Burckhardt}'s und \textsc{Damoiseau}'s Tafeln, und eine von dem Verfasser auf der Coquille gemachte chronometrische Bestimmung gab sehr übereinstimmende Resultate. — Ueber ein Mittel die Wirkung der Capillarität bei den Barometern aufzuheben, von \textsc{Laplace}. Da das Quecksilber im Gefässe nach dem Rande zu convex ist, so kann man dem die Oberfläche dieses Quecksilbers berührenden Stifte, dessen unteres Ende dem Nullpunkte der Scale entspricht, einen solchen Platz geben, dass die Wirkung der Capillarität aufgehoben wird: Herr \textsc{Brunard} hat dazu eine Hülfstafel gegeben. — Ueber die Veränderung, welche die Constante der \SourceLeaf{653} Praecession in Folge des Königsberger Katalogs der Fundamentalsterne erleidet, von \textsc{Bessel}. Der Inhalt dieses Aufsatzes ist bereits aus den astronomischen Nachrichten bekannt. --- Beobachtungen und Elemente des am 19. Mai 1825 entdeckten Cometen, von \textsc{Gambart}. --- Ueber die Anziehung der Sphäroide von \textsc{Poisson}. Dieser gehaltreiche Aufsatz besteht aus zwei Abtheilungen. Die erste beschäftigt sich mit der Entwickelung der Kugelfunctionen (so möchten wir die Functionen zweier veränderlichen Grössen, die allgemein jeden Punkt einer Kugelfläche bestimmen, nennen) in Reihen, die nach einem bekannten von den Analysten vielfach behandelten Gesetze fortschreiten, und ist mit der diesem grossen Geometer eigenthümlichen Eleganz durchgeführt. Der für das Fundamentaltheorem dieser Verwandlung schon sonst gegebene Beweis hat hier, um einigen dagegen gemachten Entwürfen zu begegnen, noch einen Zusatz erhalten, der, richtig verstanden, allerdings alle Schwierigkeiten hebt, obwohl eine vollständigere Unterscheidung der dabei vorkommenden unendlich kleinen Grössen die Evidenz des Beweises noch vollkommener machen würde; wenigstens zeigt das, was neuerlich\HistoricalNote{37} von einem berühmten englischen Geometer im Maistück des Philosophical Magazine gegen die neue \textsc{Poisson}'sche Darstellung vorgebracht ist, wenn es auch das Wesen derselben gar nicht trifft, dass diese Darstellung noch missverstanden werden konnte. Die zweite Abtheilung enthält allgemeine Untersuchungen über die Anziehung beliebiger Körper, mit vieler analytischer Kunst und tiefer als sonst bisher geschehen war, ausgeführt: doch beruht alles auf Entwickelung in Reihen, und Resultate ergeben sich nur in so fern, als die Abweichung des Körpers von der Kugelgestalt als klein vorausgesetzt wird. Für den Fall, wo diese Bedingung wegfällt, ist, nach des Verfassers Ausdruck ‘die Gestalt des Gleichgewichts von den Geometern noch nicht bestimmt, und wenn das Ellipsoid die einzige Figur wäre, wobei das Gleichgewicht möglich ist, so würde die sonderbare Folgerung nothwendig sein, dass das Gleichgewicht unmöglich wird für eine Geschwindigkeit der Rotationsbewegung, die doch noch nicht diejenige ist, bei welcher die Flüssigkeit anfängt sich zu zerstreuen.’ Wir würden uns lieber so ausdrücken, dass wir noch keinen mathematisch strengen Beweis besitzen, dass das Ellipsoid die einzig mögliche Figur für das Gleichgewicht ist: denn in der That enthält jene Folgerung durchaus nichts widersprechendes.
